\documentclass[
  oneside,
  fontset=fandol
]{ujsthesis}

% Bibliography is GB/T 7714-2015 numeric and follows citation order.
\addbibresource{bib/references.bib}

\ujstitle{面向本科毕业设计的可复现论文模板设计}
\ujstitleen{Design of a Reproducible Thesis Template for Undergraduate Projects}
\ujscollege{计算机科学与通信工程学院}
\ujsmajor{计算机科学与技术 2201 班}
\ujsstudent{张三}
\ujsstudentid{2020123456}
\ujsadvisor{李四}
\ujsrank{副教授}
\ujsdate{2026 年 6 月}
\ujsdeclarationdate{2026 年 6 月\quad 日}
\ujskeywords{本科毕业设计；LaTeX；格式规范；可复现构建}
\ujskeywordsEN{undergraduate thesis; LaTeX; formatting specification; reproducible build}
\ujsabstract{本文围绕江苏大学本科毕业设计（论文）通用格式的 LaTeX 实现问题，设计并实现一个适用于理工农医教类专业的可配置、可复现论文模板。研究首先将学校 2026 版 Word 模板中的封面、原创性声明、摘要、目录、正文、参考文献、附录和致谢等内容拆分为可维护的源文件模块，再把 A4 页面、页边距、中文和英文字体、字号、段落缩进、行距、标题层级、页眉和页码等要求集中到文档类中管理。模板采用 XeLaTeX 编译中文内容，使用随 TeX Live 提供的 Fandol 字体保证跨平台构建，并以 TeX Gyre Termes 作为英文和数字字体的可复现替代。项目同时提供图表、公式、代码、标签、交叉引用和 GB/T 7714-2015 数字型参考文献的基础接口，使论文作者能够在保持结构清晰的情况下逐步补充自己的研究内容。示例工程按照封面、前置页、正文、参考文献、致谢和附录的顺序组织，配合 latexmk、Biber 和 Makefile 完成自动重复编译。构建验证检查了 PDF 是否生成、页数和纸张规格、页眉文字、罗马数字前置页码、阿拉伯数字正文页码、目录编号、参考文献输出以及关键中文字符是否正常嵌入。结果表明，该模板能够在当前 TeX Live 环境中稳定生成结构完整的本科论文样例，并为后续替换题目、作者信息、章节内容、实验数据和图表提供统一的排版入口。}
\ujsabstractEN{This project develops a configurable and reproducible LaTeX template for undergraduate design reports and theses in science, engineering, agriculture, medicine, and education at Jiangsu University. The work begins by translating the 2026 Word template into maintainable source modules for the cover, originality statement, Chinese abstract, English abstract, table of contents, main chapters, references, acknowledgements, and appendix. Page geometry, CJK and Latin fonts, font sizes, paragraph indentation, line spacing, heading levels, headers, and page numbering are centralized in one document class so that content files do not need to repeat formatting instructions. XeLaTeX is used for Chinese typesetting, Fandol fonts supplied by TeX Live provide a portable CJK environment, and TeX Gyre Termes provides a reproducible Latin font substitute for Times New Roman. The project also exposes basic interfaces for figures, tables, equations, source code, labels, cross references, and numeric references formatted according to GB/T 7714-2015. The source tree follows a stable writing order: metadata and bibliography are configured in the root file, chapters are maintained independently, and the class controls the visual rules shared by all pages. A latexmk configuration, a Makefile, and a GitHub Actions workflow provide repeatable local and continuous integration builds. The example build checks PDF generation, A4 page geometry, required header text, Roman numbering in the front matter, Arabic numbering from the introduction through the appendix, Arabic chapter numbers in the contents, bibliography generation by Biber, and embedding of the selected Chinese and Latin fonts. It also keeps the layout logic separate from the prose so that a change to a title or a paragraph does not require manual editing of page headers or table of contents entries. The modular design makes it possible to add a discipline-specific chapter, insert a measured experiment, or replace the bibliography database without changing the document class. The resulting sample demonstrates a complete compilation path from editable source to PDF. It is intended as a practical starting point: authors still need to replace the placeholder metadata and example prose, add their own figures and measurements, verify the latest college-level requirements, and inspect the final PDF on the target printer before submission.}

\begin{document}

\MakeCover
\UJSFrontMatter
\MakeDeclaration

\MakeAbstract
\MakeTOC

\UJSMainMatter
\chapter{绪论}

\section{课题背景和意义}
本科毕业设计既是对专业知识的综合运用，也是对研究过程和工程表达能力的集中检验。一个可复现的论文模板能够把学校的格式要求落实为源文件中的稳定规则，减少手工调整带来的版式漂移。

\subsection{研究背景}
传统文字处理模板通常依赖本地软件状态。字体、目录域、交叉引用和页码更新如果没有被持续检查，提交前容易出现不一致。LaTeX 将结构和排版分开，使章节、公式、图表和参考文献可以由源文件重建\cite{lamport1994}。

\subsection{研究意义}
本项目提供一个面向江苏大学理工农医教类毕业设计（论文）的起点工程。使用者只需要修改元数据和章节文件，就可以生成与学校模板约束相符的初稿，并把注意力放在研究内容、数据和结论上。

\section{国内外相关工作}
学术论文模板通常包含文档类、字体配置、前置内容、章节组织、参考文献样式和构建脚本。相关工作表明，模板应同时提供清晰的默认值和可覆盖的接口，避免使用者直接修改类文件\cite{mittelbach2004}。本文以模块化源文件和自动化构建为核心，保留后续加入学科专用命令的空间。

\section{本文主要工作}
本文的主要工作包括：提取学校 2026 版通用格式中的可执行约束；实现 `ujsthesis` 文档类；组织封面、声明、摘要、目录、正文、参考文献、附录和致谢；最后通过 XeLaTeX、Biber 和 PDF 文本检查验证样例工程。

\chapter{模板设计与实现}

\section{总体结构}
项目采用与参考项目相近的“入口文件—章节文件—样式文件—构建脚本”结构，但文档类按照江苏大学模板独立实现。入口文件只负责元数据、文献数据库和章节顺序；格式细节集中在 \texttt{style/ujsthesis.cls} 中。

\begin{table}[H]
  \centering
  \caption{项目文件与职责}
  \begin{tabular}{p{0.30\linewidth}p{0.56\linewidth}}
    \toprule
    文件或目录 & 职责 \\
    \midrule
    \texttt{main.tex} & 配置元数据并组织编译顺序 \\
    \texttt{style/ujsthesis.cls} & 实现字体、尺寸、页眉页码和标题层级 \\
    \texttt{chapters/} & 保存可独立维护的论文内容 \\
    \texttt{bib/references.bib} & 使用 GB/T 7714-2015 的参考文献数据 \\
    \texttt{Makefile}、\texttt{latexmkrc} & 提供可复现的本地构建入口 \\
    \bottomrule
  \end{tabular}
\end{table}

\section{字体与页面设置}
正文使用小四号中文字体和 1.5 倍近似行距，英文和数字使用 TeX Gyre Termes 作为 Times New Roman 的跨平台替代。页面使用 A4 纸，页边距根据所附 2026 版 Word 模板转换结果设置；如学院另发细则，应优先以学院细则为准并在类文件中调整。

\section{图表和公式接口}
图表、公式和代码应在正文中被引用。公式示例为
\begin{equation}
  y = f(x;\theta) + \varepsilon,
  \label{eq:model}
\end{equation}
其中 $\theta$ 是待估参数，$\varepsilon$ 表示未建模误差。式~\eqref{eq:model} 只示范交叉引用，不代替具体课题的模型假设。

\chapter{构建验证与使用示例}

\section{构建流程}
在安装 TeX Live、XeLaTeX 和 Biber 后，可在项目根目录执行：
\begin{lstlisting}[language=bash,caption={本地构建命令}]
make
make clean
\end{lstlisting}
构建脚本会运行 XeLaTeX、Biber 和必要的重复编译，使目录、引用和参考文献稳定下来。

\section{验证范围}
样例工程覆盖封面、原创性声明、中英文摘要、目录、四个正文级章节、参考文献、附录和致谢。自动检查包括 PDF 是否生成、页数是否非零、关键页眉文字是否出现、中文和英文摘要标题是否存在，以及引用是否能由 Biber 解析。上述检查只能证明样例工程在当前 TeX Live 环境中可构建，不能替代学院对最终论文内容、页数和打印样张的审核。

\section{使用注意事项}
中文摘要应达到学校要求的 400 字以上，英文摘要应达到 350 个实词以上；参考文献应按 GB/T 7714-2015 书写，正文引用顺序与文后编号应保持一致。提交前应在目标打印机或 PDF 阅读器中检查封面、页眉、图表和分页，并根据导师或学院的最新通知更新字段。

\chapter{结论}

\section{主要结论}
本项目将江苏大学 2026 版通用格式中的主要版式要求转化为独立的 LaTeX 文档类和示例工程，形成了从元数据配置到 PDF 输出的连续路径。章节文件、文献数据库和构建脚本分离后，论文内容可以在不修改格式实现的情况下持续迭代。

\section{不足与展望}
当前模板未内置学校官方校徽图像，也未替代不同学院可能附加的任务书、评阅表和装订要求。后续可以在取得学院授权的前提下加入官方封面图形，增加学院专用选项，并将样例扩展为真实课题的图表和实验数据。


\UJSReferences

\UJSBackMatterHeading{致谢}
\ujsbodyfont
感谢指导教师在选题、结构设计和论文修改过程中提供的指导，也感谢同学们在格式测试和构建验证中给予的帮助。本文的示例内容仅用于说明模板的使用方式，正式论文应替换为作者本人完成且经过核验的研究内容。


\UJSBackMatterHeading{附录}
\appendix
\chapter{构建环境记录}

本附录记录样例构建所需的最小环境：TeX Live 2026、XeLaTeX、Biber、GNU Make 和常见字体。环境版本会随发行版更新，实际提交时应记录本机的 \texttt{xelatex --version} 和 \texttt{biber --version} 输出。

\section{最小命令}
\begin{lstlisting}[language=bash,caption={环境检查命令}]
xelatex --version
biber --version
make
\end{lstlisting}


\end{document}
